\subsection{Registro conjunto de un resultado y su plazo configurado}\label{sec:impl-resultado-plazo-derivado}

El registro de una sesión declarada y el alta de un plazo pueden formar una sola instrucción explícita del operador. El servicio prepara un resultado ordinario de audiencia y un plazo, ambos en su revisión inicial, dentro del mismo expediente. La selección temporal del plazo referencia exactamente la audiencia, la identidad del resultado y su revisión uno; si selecciona un acuerdo, éste debe existir en los valores propuestos. Las identidades y las dos operaciones permanecen estables durante revisión, confirmación y recuperación. El registro ordinario de resultados sin plazo conserva su propio recorrido. Una consecuencia por instrucción limita esta interfaz, sin afirmar que una actuación sólo pueda producir un plazo.

La composición utiliza el perfil y el evaluador descritos en las \cref{sec:implementacion-perfiles-plazos,sec:implementacion-plazos-persistentes}. Exige seleccionar perfil, responsable y políticas de seguimiento; el calendario, la cantidad ordenada, las condiciones y la calificación temporal se declaran cuando corresponde. La selección histórica de una dependencia se conserva separada de su cabeza observada. El relato, la comparecencia o el estado de sesión concluida no determinan una hora de finalización, una notificación eficaz ni la regla jurídica aplicable. Los datos ausentes pueden producir un bloqueo explicado o una candidata civil; una referencia ajena o una evidencia inconsistente se rechazan en lugar de presentarse como bloqueo aceptable.

\subsubsection{Revisión prospectiva y compromisos separados}

Durante la preparación todavía no existe la revisión del resultado que servirá como fuente persistida. La aplicación resuelve y verifica las fuentes existentes y calcula sobre los valores propuestos mediante una vía prospectiva específica. No construye un detalle histórico ficticio para satisfacer el contrato ordinario de insumos. El reloj observado permite comprobar la declaración, pero no se convierte en su instante definitivo de registro. Tampoco se fabrican un evento de fuente, un recibo final del plazo ni una captura administrativa nueva.

La representación canónica \texttt{HRDL1} vincula el comando conjunto, autor, programación exacta, continuación opcional, asistentes, soporte, administración, perfil, calendario, responsable, políticas y cálculo revisado. Su huella representa la instrucción presentada para aprobación. Una modificación exige preparar y revisar otra vez, incluso cuando conserve la misma fecha calculada. Los formatos \texttt{HRES1} y \texttt{HRTX1} mantienen su significado; el resultado ordinario no adquiere una huella de captura propia por participar en la composición.

La confirmación produce un compromiso distinto, \texttt{HRDC1}, que liga la revisión prospectiva con los recibos reales del resultado y del plazo, el evento exacto y el instante de captura. La huella revisada y la de captura no son intercambiables. Ambas permiten comprobar la correspondencia entre lo aprobado y lo almacenado, pero no constituyen firma judicial, sello de tiempo ni acreditación jurídica del acto.

\subsubsection{Autorización y escritura atómica}

El servicio exige simultáneamente las capacidades de gestionar resultados y plazos: Owner, o Litigante asignado al expediente. Paralegal y Cliente no gestionan esta operación. Se comprueba el principal autenticado antes de preparar y después de resolver el material; antes de confirmar y de entregar el registro se vuelve a autenticar. La admisión del soporte ordinario descifra y valida la versión exacta fuera de la transacción que obtuvo las referencias. La preparación no reserva identificadores ni concede autoridad para una escritura posterior.

El adaptador PostgreSQL obtiene el bloqueo común de auditoría, reautoriza y vuelve a resolver el material revisado. Compara también los registros cifrados que sustentaron la admisión. Dentro de la misma transacción inserta el resultado ordinario R1, obtiene su evento real y prepara el plazo mediante la vía ordinaria que ya puede consultar esa fuente persistida. El cálculo final debe corresponder a la revisión aprobada. El reloj definitivo del resultado se conserva en la captura conjunta; una falla en el plazo, su origen o la auditoría revierte todas las escrituras.

La familia de migraciones \texttt{0032} incorpora el origen inmutable que relaciona operaciones, resultado y plazo R1, autor original, evento, compromisos canónicos y auditoría. No añade una restricción global que prohíba otros plazos legítimos derivados de la misma fuente. Utiliza el evento que ya emite el resultado, sin duplicarlo, elegir el mayor número de secuencia ni marcarlo procesado para simular la ejecución del consumidor. Los huecos de secuencia causados por transacciones revertidas son admisibles. El arranque valida catálogo, permisos e inventario histórico antes de admitir el uso del almacén.

\subsubsection{Historia exacta y conciliación del intento}

La reconstrucción histórica verifica ambos componentes, sus recibos, fuentes y cabezas capturadas, el evento y el origen auditado. Reproduce los bytes del compromiso a partir del cálculo conservado; no vuelve a ejecutar la aritmética ni sustituye las dependencias por sus cabezas actuales. El correo y rol capturados del autor siguen siendo evidencia histórica, mientras el acceso se decide mediante la identidad y autoridad actuales. La mera existencia de dos registros independientes con valores semejantes no permite adoptarlos como una creación conjunta.

Ante una respuesta perdida, la consulta explícita del mismo comando puede devolver la creación original. La aplicación contrasta sus compromisos con aquel intento y conserva las dos revisiones iniciales. Una confirmación idéntica concurrente se serializa bajo el bloqueo compartido y recupera el mismo origen; una instrucción distinta entra en conflicto. La conciliación no crea automáticamente otra operación ni convierte una revisión posterior en prueba del envío anterior.

\subsubsection{Transporte y recuperación del editor}

{\raggedright Las rutas de preparación y confirmación se sitúan bajo el resultado de la audiencia y terminan en \texttt{derived-deadline/prepare} y \texttt{derived-deadline/submit}. El contrato completo se conserva en \rutarepo{docs/hearing-derived-deadlines.md}. Exigen bearer, objetos JSON estrictos, ausencia de parámetros de consulta y un cuerpo máximo de 1~MiB. La preparación devuelve una instrucción revisable o un registro original recuperado; la confirmación devuelve los dos detalles R1 y el origen conjunto. La secuencia del evento se transmite como cadena decimal para evitar redondeo en JavaScript.\par}

Qadra integra el editor en la vista de resultados y reutiliza los campos ordinarios de resultado y plazo. Expone las declaraciones, políticas, fuentes, cálculo y bloqueos antes de pedir confirmación. Una preparación válida puede contener un plazo bloqueado: no expresa aprobación jurídica ni vencimiento operativo. Los borradores incompletos se conservan sin convertir silenciosamente sus campos vacíos en declaraciones de desconocimiento.

Si caduca la sesión, la reentrada del mismo usuario reautoriza el expediente y recupera el borrador sin restaurar la aprobación anterior. Un intento incierto conserva el comando y sus identidades; primero consulta su registro original, antes de cargar catálogos o soportes actuales. Una respuesta de preparación idéntica mantiene la incertidumbre y puede permitir un reenvío expreso; sólo la correspondencia con el registro conjunto confirma la captura. El cierre administrativo impide una nueva escritura, pero mantiene la conciliación autorizada. Los cambios de identidad y las denegaciones retiran los datos correspondientes, y las respuestas tardías no reactivan un editor descartado.

El alta no evita las comprobaciones posteriores de vigencia, aceptación y permiso de Agenda y alertas. Un cálculo histórico puede permanecer disponible sin fecha operativa, y un registro bloqueado no se convierte en vencimiento por haberse almacenado. El navegador con servicios reales comprobó el recorrido del editor, Agenda y avisos internos del plazo. La campaña completa de API y restauración conservó los pares iniciales calculable y bloqueado, sus orígenes, recibos y auditoría; la conciliación recuperó las creaciones originales sin duplicarlas. Estas comprobaciones tienen alcances distintos y no completan el corpus jurídico ni la validación con personal del despacho. El cierre de integración continua, la integración remota y el despliegue siguen pendientes. Las \cref{sec:pruebas-resultado-plazo-derivado,sec:matriz-resultado-plazo-derivado} separan sus fronteras de verificación.
